\begin{abstract}
Clinical image-analysis systems operate under strict limits on memory, latency, and the transmission of sensitive patient data. State-of-the-art models, however, continue to grow in size and compute, a trend that fits resource-constrained deployment poorly.
This study systematically evaluates deep learning models with fewer than one million parameters on small-to-medium-sized datasets for image classification and detection, focusing on efficiency and cross-dataset generalization. Results show that on saturated classification benchmarks, compact models can achieve performance statistically comparable to models an order of magnitude larger. The best model attains 98.3\% accuracy at 0.36 GFLOPs, running six times faster with less than a one-percentage-point accuracy difference. However, a performance gap remains on more challenging cross-dataset tasks, with errors concentrated in low-confidence predictions sensitive to input perturbations. The most efficient architecture, which uses parallel atrous convolution branches to capture multi-scale features, is particularly vulnerable to scale-sensitive inputs.
To address this limitation, we propose \textbf{MCG-TTA}, an architecture-informed inference strategy designed to improve robustness while limiting computational overhead. MCG-TTA improves slice-level F1 by up to 1.8 percentage points at approximately 1.9$\times$ the baseline latency, while retaining a model size of 0.91 million parameters. These findings suggest that, for resource-constrained clinical applications, uncertainty-aware inference can recover residual accuracy more efficiently than increasing model complexity.
\end{abstract}
\begin{IEEEkeywords}
Brain MRI, convolutional neural networks, compact neural networks, test-time augmentation, glioma grading, edge inference
\end{IEEEkeywords}
